\section*{Hoofdstuk \ref{ch:b3:representations} --- Representaties
van eindige groepen}

\begin{solution}{exo:b3:representations:1}
(a) $\Z/n\Z$ is abels: alle \omterm{def:b3:representations:rep}{irreducibele representaties} hebben
graad $1$ (\cref{prop:b3:representations:onedim}), dat wil zeggen
het zijn morfismen $\chi \colon \Z/n\Z \to \C^\times$, vastgelegd
door $\chi(\bar 1) = \omega$ met $\omega^n = 1$: de $n$ \omterm{def:b3:representations:character}{karakters}
$\chi_k(\bar m) = \eu^{2\iu\pi km/n}$. Voor $n = 4$ (de klassen
zijn de elementen $\bar0,\bar1,\bar2,\bar3$):
\[
\begin{array}{c|cccc}
& \bar0 & \bar1 & \bar2 & \bar3\\
\hline
\chi_0 & 1 & 1 & 1 & 1\\
\chi_1 & 1 & \iu & -1 & -\iu\\
\chi_2 & 1 & -1 & 1 & -1\\
\chi_3 & 1 & -\iu & -1 & \iu
\end{array}
\]
(b) Rijen: $\langle\chi_k,\chi_l\rangle = \frac14\sum_m
\eu^{2\iu\pi(l-k)m/4} = \delta_{kl}$ (meetkundige som); de
kolommen net zo. De matrix $(\eu^{2\iu\pi km/n})_{k,m}$ is de
matrix van de \emph{discrete Fouriertransformatie} --- dezelfde
eenheidswortelfilter die in het hoofdstuk over genererende
functies in het volume van bachelorjaar 2 werd gebruikt; de
\omterm{thm:b3:representations:orthogonality}{orthogonaliteit van karakters} veralgemeent de inversieformule van
die transformatie.
\end{solution}

\begin{solution}{exo:b3:representations:2}
(a) De matrix van $\rho(g)$ in de basis $(e_x)$ is een
permutatiematrix, met als spoor het aantal $x$ met $g\cdot x = x$.
Dan is
\[
\langle\mathbf 1, \chi\rangle = \frac1{\abs G}\sum_g
\abs{\operatorname{Fix}(g)} = \#\{\text{banen}\}
\]
volgens de teltelling van Burnside (\cref{exo:b3:groups:5}) ---
gelijkwaardig berekent dit de multipliciteit van de triviale
\omterm{def:b3:representations:rep}{representatie}, waarvan de isotypische ruimte de ruimte van
$G$-invariante vectoren is, van dimensie het aantal \omterm{def:b3:groups:action}{banen} (één
indicatorfunctie per \omterm{def:b3:groups:action}{baan}).

(b) Omdat $\operatorname{Fix}_{X\times X}(g) =
\operatorname{Fix}_X(g)^2$, geeft onderdeel (a) toegepast op $X
\times X$ dat $\langle\chi,\chi\rangle = \frac1{\abs
G}\sum\abs{\operatorname{Fix}(g)}^2 = \#\{\text{banen op }
X\times X\}$ (want $\chi$ is reëel). Schrijven we $\chi = \mathbf
1 + \psi$, dan is $\langle\mathbf1,\chi\rangle = 1$
(transitiviteit), dus $\langle\psi,\psi\rangle =
\langle\chi,\chi\rangle - 1$. De \omterm{def:b3:groups:action}{werking} op $X\times X$ heeft de
diagonaal als één \omterm{def:b3:groups:action}{baan}; er is precies één andere \omterm{def:b3:groups:action}{baan} dan en
slechts dan als $G$ transitief is op de paren verschillende
punten: $\langle\psi,\psi\rangle = 1$ precies wanneer de \omterm{def:b3:groups:action}{werking}
$2$-transitief is (\cref{cor:b3:representations:multiplicity}).

(c) $S_n$ is $2$-transitief op $\{1,\dots,n\}$ (stuur een
willekeurig paar verschillende punten waarheen je wilt): $\psi =
\chi_{\mathrm{std}}$ is \omterm{def:b3:representations:rep}{irreducibel}.
\end{solution}

\begin{solution}{exo:b3:representations:3}
$S_3$: drie klassen, $\sum n_i^2 = 6 = 1 + 1 + 4$; de twee
\omterm{def:b3:representations:character}{karakters} van graad $1$ zijn $\mathbf 1$ en $\varepsilon$ (want
$[S_3 : A_3] = 2$); de derde rij $(2, a, b)$ volgt uit de
\omterm{thm:b3:representations:orthogonality}{kolomorthogonaliteit} met de kolom van $e$: $1 - 1 + 2a = 0$ en $1
+ 1 + 2b = 0$, dus $a = 0$ en $b = -1$ --- de tafel van
\cref{ex:b3:representations:s3}.

Controles voor $S_4$ (rijen): $\langle\chi_{\mathrm{std}},
\varepsilon\chi_{\mathrm{std}}\rangle = \frac1{24}(9 - 6 + 3 + 0 -
6) = 0$ en $\langle\chi_2,\chi_2\rangle = \frac1{24}(4 + 0 + 12 +
8 + 0) = 1$. Kolommen: $e$ tegen $(1\,2)$: $1 - 1 + 0 + 3 - 3 =
0$; en $(1\,2)$ tegen zichzelf: $1 + 1 + 0 + 1 + 1 = 4 =
\abs{Z_{S_4}((1\,2))} = 24/6$.

Permutatiekarakter op $4$ punten: $(4, 2, 0, 1, 0) = \mathbf 1 +
\chi_{\mathrm{std}}$ (de aantallen vaste punten; trek de bovenste
rij af). Voor $\chi_{\mathrm{std}}^2 = (9, 1, 1, 0, 1)$:
\[
\langle\mathbf1,\cdot\rangle = \tfrac{9 + 6 + 3 + 0 + 6}{24} =
1,\quad
\langle\varepsilon,\cdot\rangle = 0,\quad
\langle\chi_2,\cdot\rangle = \tfrac{18 + 6}{24} = 1,\quad
\langle\chi_{\mathrm{std}},\cdot\rangle = 1,\quad
\langle\varepsilon\chi_{\mathrm{std}},\cdot\rangle = 1:
\]
dus $\chi_{\mathrm{std}}^2 = \mathbf 1 + \chi_2 +
\chi_{\mathrm{std}} + \varepsilon\chi_{\mathrm{std}}$ (dimensies:
$9 = 1 + 2 + 3 + 3$).
\end{solution}

\begin{solution}{exo:b3:representations:4}
Beide groepen hebben vijf klassen en het gradenpatroon
$(1,1,1,1,2)$ (vier \omterm{def:b3:representations:character}{karakters} van graad $1$ uit de abelianisering
$\cong (\Z/2\Z)^2$, en dan $\sum n_i^2 = 8$). Ordenen we de
klassen als $e$, $z$ (de centrale involutie: $r^2$ respectievelijk
$-1$) en de drie klassen met twee elementen:
\[
\begin{array}{c|ccccc}
& e & z & C_1 & C_2 & C_3\\
\hline
\chi^{(1)} & 1 & 1 & 1 & 1 & 1\\
\chi^{(2)} & 1 & 1 & 1 & -1 & -1\\
\chi^{(3)} & 1 & 1 & -1 & 1 & -1\\
\chi^{(4)} & 1 & 1 & -1 & -1 & 1\\
\chi^{(5)} & 2 & -2 & 0 & 0 & 0
\end{array}
\]
(de laatste rij uit de \omterm{thm:b3:representations:orthogonality}{kolomorthogonaliteit}). Identieke tafels
voor $D_4$ en $Q_8$, die niet isomorf zijn
(\cref{pb:b3:groups:1}). De tafel legt \emph{wel} vast: $\abs G$,
de klassegroottes, het \omterm{ex:b3:groups:actions}{centrum} ($\{g : \abs{\chi_i(g)} = n_i\
\forall i\}$: orde $2$ in beide), de abelianisering en het hele
tralie van \omterm{def:b3:groups:normal}{normaaldelers} (kernen en doorsneden,
\cref{exo:b3:representations:6}). Ze legt de ordes van de
elementen \emph{niet} vast: $D_4$ heeft vijf involuties en $Q_8$
maar één --- het isomorfietype is dus werkelijk fijner dan de
\omterm{def:b3:representations:table}{karaktertafel}.
\end{solution}

\begin{solution}{exo:b3:representations:5}
(a) In $S_4$ heeft de centralisator van $(1\,2\,3)$ orde $24/8 =
3$: het is $\langle(1\,2\,3)\rangle \subseteq A_4$. Dus heeft
$Z_{A_4}((1\,2\,3))$ orde $3$ en telt de $A_4$-klasse $12/3 = 4$
elementen: de acht $3$-cykels vallen uiteen in twee
$A_4$-klassen (met representanten $(1\,2\,3)$ en haar inverse).

(b) $A_4/V \cong \Z/3\Z$ geeft drie \omterm{def:b3:representations:character}{karakters} van graad $1$ (met
$\omega = \eu^{2\iu\pi/3}$; de klassen van $3$-cykels gaan naar
$\bar1$ en $\bar2$); de beperking van $\chi_{\mathrm{std}}$ blijft
\omterm{def:b3:representations:rep}{irreducibel} ($\langle\chi,\chi\rangle = \frac1{12}(9 + 3 \cdot 1 +
0 + 0) = 1$):
\[
\begin{array}{c|cccc}
A_4 & e\,[1] & (1\,2)(3\,4)\,[3] & (1\,2\,3)\,[4] &
(1\,3\,2)\,[4]\\
\hline
\mathbf 1 & 1 & 1 & 1 & 1\\
\chi_\omega & 1 & 1 & \omega & \omega^2\\
\chi_{\bar\omega} & 1 & 1 & \omega^2 & \omega\\
\chi_3 & 3 & -1 & 0 & 0
\end{array}
\]
(c) Kernen: $\ker\chi_\omega = \ker\chi_{\bar\omega} = V$ en
$\ker\chi_3 = \{e\}$ (geen enkele andere ingang heeft modulus
$3$). De \omterm{def:b3:groups:normal}{normaaldelers} zijn de doorsneden van de kernen
(\cref{exo:b3:representations:6}): $\{e\}$, $V$ en $A_4$ --- in
het bijzonder heeft $A_4$ geen \omterm{def:b3:groups:normal}{normaaldeler} van orde $2$ of van
index $2$.
\end{solution}

\begin{solution}{exo:b3:representations:6}
(a) Is $\chi(g) = \chi(e) = n$, dan dwingt de gelijkheid in
$\abs{\chi(g)} \leq n$ af dat $\rho(g) = \lambda\,\mathrm{id}$
(\cref{prop:b3:representations:charbasics}) met $n\lambda = n$:
dus $\rho(g) = \mathrm{id}$. De omkering is duidelijk.

(b) Zij $N \trianglelefteq G$. De \omterm{ex:b3:representations:examples}{reguliere representatie} van
$G/N$ is trouw; ontbind haar in \omterm{def:b3:representations:rep}{irreducibele representaties} van
$G/N$ en til die op naar $G$
(\cref{prop:b3:representations:onedim}(b)): dat geeft irreducibele
\omterm{def:b3:representations:character}{karakters} $\chi_{i_1}, \dots$ van $G$ waarvan de kernen $N$
bevatten en waarvan de \emph{gemeenschappelijke} kern precies het
origineel van $\{e\}$ is, dat wil zeggen $N$ (trouwheid op het
quotiënt). Dus $N = \bigcap_j \ker\chi_{i_j}$.

(c) Is $G$ enkelvoudig, dan is voor een niet-triviaal \omterm{def:b3:representations:rep}{irreducibel}
$\chi$ de \omterm{def:b3:groups:normal}{normaaldeler} $\ker\chi$ niet gelijk aan $G$ (een
\omterm{def:b3:representations:rep}{irreducibele representatie} die op heel $G$ triviaal is, is het
triviale \omterm{def:b3:representations:character}{karakter}), dus $\ker\chi = \{e\}$. Stel omgekeerd dat
alle niet-triviale kernen triviaal zijn, en zij $N
\trianglelefteq G$ met $N \neq G$. In de uitdrukking van $N$ als
doorsnede van kernen uit (b) is een van de betrokken \omterm{def:b3:representations:character}{karakters}
niet-triviaal (waren ze alle triviaal, dan was de doorsnede $G$),
en de kern daarvan is $\{e\}$: dus $N = \{e\}$. De enige
\omterm{def:b3:groups:normal}{normaaldelers} zijn dus $\{e\}$ en $G$.
\end{solution}

\begin{solution}{exo:b3:representations:7}
Orde $8$, niet-abels: het aantal \omterm{def:b3:representations:character}{karakters} van graad $1$ is $[G :
D(G)]$, een echte deler van $8$ (niet-abels: $D(G) \neq \{e\}$),
en $\sum n_i^2 = 8$. Met $k$ enen en de overige graden $\geq 2$
moet $8 - k$ een som van kwadraten $\geq 4$ zijn, met $k \mid 8$
en $k < 8$. Voor $k = 4$: één graad $2$ --- consistent. Voor $k =
2$: er blijft $6$ over, geen som van kwadraten $\geq 4$. Voor $k =
1$: onmogelijk, want $k = [G : D(G)] \geq 2$ --- $G/D(G)$ is een
niet-triviale abelse $2$-groep, aangezien $G$ een $2$-groep is met
$D(G) \neq G$ wegens de \omterm{def:b3:groups:derived}{oplosbaarheid} van $p$-groepen
(\cref{ex:b3:groups:solvableexamples}). Het patroon is dus
$(1,1,1,1,2)$. Trouwheid van $\chi_5$: de vier \omterm{def:b3:representations:character}{karakters} van graad
$1$ bevatten alle $D(G)$ in hun kern; lag er een minimale
niet-triviale \omterm{def:b3:groups:normal}{normaaldeler} $N \subseteq \ker\chi_5$, dan lag $N$
in alle vijf kernen, waarvan de doorsnede triviaal is (de
\omterm{ex:b3:representations:examples}{reguliere representatie} is trouw): tegenspraak. Orde $p^3$,
niet-abels: $Z(G)$ heeft orde $p$ (bij orde $p^2$ zou $G/Z$
cyclisch en $G$ abels zijn), $G/Z(G)$ van orde $p^2$ is abels, dus
$D(G) \subseteq Z(G)$, en $D(G) \ne \{e\}$: dus $D(G) = Z(G)$, wat
$p^2$ \omterm{def:b3:representations:character}{karakters} van graad $1$ geeft. De overige graden voldoen aan
$\sum n_i^2 = p^3 - p^2$ met elke $n_i > 1$ die $\abs G$ deelt
(\cref{pb:b3:representations:1}, vraag 8), dus $n_i = p$ ($n_i =
p^2$ zou te groot zijn: $p^4 > p^3 - p^2$): precies $p - 1$
\omterm{def:b3:representations:character}{karakters} van graad $p$.
\end{solution}

\begin{solution}{exo:b3:representations:8}
Definieer voor \omterm{def:b3:representations:rep}{representaties} $\rho, \sigma$ van $G$ en $H$ op $V$
en $W$ de \omterm{def:b3:representations:rep}{representatie} $\rho \boxtimes \sigma$ van $G \times H$
op $V \otimes W$ --- concreet, op matrices: $(\rho\boxtimes
\sigma)(g,h)$ is het kroneckerproduct $\rho(g)\otimes \sigma(h)$,
met spoor $\operatorname{tr}\rho(g)\operatorname{tr}\sigma(h) =
\chi(g)\psi(h)$ (het kroneckerproduct $A \otimes B$ heeft spoor
$\operatorname{tr}A\operatorname{tr}B$: de diagonaal bestaat uit
de $a_{kk}b_{ll}$). Dus is $\chi\psi$ een \omterm{def:b3:representations:character}{karakter}, en
\[
\langle\chi\psi, \chi'\psi'\rangle_{G\times H}
= \frac{1}{\abs G\abs H}\sum_{g,h}
\overline{\chi(g)\psi(h)}\,\chi'(g)\psi'(h)
= \langle\chi,\chi'\rangle_G\,\langle\psi,\psi'\rangle_H
= \delta_{\chi\chi'}\delta_{\psi\psi'}.
\]
In het bijzonder is $\langle\chi\psi,\chi\psi\rangle = 1$: elke
$\chi\psi$ is \omterm{def:b3:representations:rep}{irreducibel}
(\cref{cor:b3:representations:multiplicity}). Dat zijn $r_Gr_H$
verschillende irreducibele \omterm{def:b3:representations:character}{karakters}; en de klassen van $G\times
H$ zijn de producten van klassen ($(g,h) \sim (g',h')$ per
component), dus zijn er $r_Gr_H$: de lijst is volledig
(\cref{thm:b3:representations:numberirr}). Voor $(\Z/2\Z)^2$: de
vier tekenkarakters $(\pm1)\otimes(\pm1)$ --- het linksboven
gelegen blok van de tafel uit \cref{exo:b3:representations:4}. Een
eindige abelse groep is een product van cyclische groepen
(\cref{cor:b3:modules:abelian}); haar irreducibele \omterm{def:b3:representations:character}{karakters} zijn
de producten van de cyclische \omterm{def:b3:representations:character}{karakters}
(\cref{exo:b3:representations:1}): alle van graad $1$.
\end{solution}

\begin{solution}{exo:b3:representations:9}
(a) $D(A_5) = A_5$ (enkelvoudig en niet-abels), dus is het enige
\omterm{def:b3:representations:character}{karakter} van graad $1$ het triviale
(\cref{prop:b3:representations:onedim}). We hebben $n_2^2 + n_3^2
+ n_4^2 + n_5^2 = 59$ nodig met elke $n_i \geq 2$; met de
beschikbare kwadraten $4, 9, 16, 25, 36, 49$ is de enige
multiverzameling die werkt $\{9, 9, 16, 25\}$: met grootste
kwadraat $49$ blijft $10$ over, geen som van drie kwadraten $\geq
4$; met $36$ blijft $23$ over, ook niet ($16 + 4 + 4 = 24$, $9 + 9
+ 4 = 22$); met grootste $25$ gaat $25 + 16 + 9 + 9 = 59$ wél op
en falen de varianten $25 + 25$, $25 + 16 + 16$ en $25 + 16 + 4$;
en met grootste $16$: $16\cdot3 = 48 < 59 - 4$. Graden: $1, 3, 3,
4, 5$.

(b) Permutatie op $5$ punten: vaste punten $(5, 1, 2, 0, 0)$, dus
$\chi_4 = (4, 0, 1, -1, -1)$ met $\langle\chi_4,\chi_4\rangle =
\frac{16 + 0 + 20 + 12 + 12}{60} = 1$: \omterm{def:b3:representations:rep}{irreducibel}. \omterm{def:b3:groups:action}{Werking} op de
zes Sylow-$5$-deelgroepen: een involutie houdt er precies $2$ vast
(de \omterm{ex:b3:groups:actions}{normalisatoren} zijn diëdraal van orde $10$ en bevatten elk $5$
involuties: $30$ incidenties voor $15$ involuties), een element
van orde $3$ houdt er $0$ vast (in $D_5$ zit geen orde $3$) en een
element van orde $5$ precies $1$ (het ligt in een unieke
\omterm{def:b3:groups:sylow}{Sylow-deelgroep}): permutatiekarakter $(6, 2, 0, 1, 1)$, en $\chi_5
= (5, 1, -1, 0, 0)$ met norm $\frac{25 + 15 + 20 + 0 + 0}{60} =
1$: \omterm{def:b3:representations:rep}{irreducibel}. Er blijven twee rijen $(3, a, b, c, d)$ en $(3,
a', b', c', d')$ over. Kolomnormen
(\cref{cor:b3:representations:column}): op de klasse van
$(1\,2)(3\,4)$ is $\abs{Z} = 4$, dus $1 + 0 + 1 + a^2 + a'^2 = 4$;
en die kolom tegen $e$: $1\cdot1 + 4\cdot 0 + 5\cdot1 + 3(a + a')
= 0$, dus $a + a' = -2$ en $a^2 + a'^2 = 2$: $a = a' = -1$. Op de
$3$-cykels is $\abs Z = 3$: $1 + 1 + 1 + b^2 + b'^2 = 3$, dus $b =
b' = 0$. Op elke $5$-klasse is $\abs Z = 5$: de kolom tegen $e$
luidt $1\cdot 1 + 4\cdot(-1) + 5\cdot 0 + 3(c + c') = 0$, dus $c +
c' = 1$; en $c^2 + c'^2 = 3$: $\{c, c'\} =
\bigl\{\frac{1+\sqrt5}2, \frac{1-\sqrt5}2\bigr\}$ --- de gulden
snede en haar toegevoegde; de tweede $5$-klasse draagt de
verwisselde waarden (de twee rijen moeten orthogonaal zijn).

(c) In de afgemaakte tafel is buiten de eerste kolom geen enkele
ingang van een niet-triviale rij gelijk aan haar graad: elke kern
$\{g : \chi_i(g) = n_i\}$ is triviaal. Volgens
\cref{exo:b3:representations:6}(c) is $A_5$ dus enkelvoudig.
\end{solution}

\begin{solution}{exo:b3:representations:10}
$\rho(z)$ commuteert met elke $\rho(g)$ ($z$ is centraal), dat wil
zeggen $\rho(z) \in \operatorname{End}_G(V) = \C\,\mathrm{id}$
(Schur): dus $\rho(z) = \lambda_z\,\mathrm{id}$, en $z \mapsto
\lambda_z$ is multiplicatief, dus een morfisme $Z(G) \to
\C^\times$. Dan is $\chi(z) = n\lambda_z$ met $\abs{\lambda_z} =
1$ (een eenheidswortel): $\abs{\chi(z)} = n$. Is $\rho$ trouw, dan
is $\lambda$ injectief op $Z(G)$ ($\rho(z) = \mathrm{id} \iff
\lambda_z = 1$), zodat $Z(G)$ zich in $\C^\times$ inbedt; en een
eindige deelgroep van de multiplicatieve groep van een lichaam is
cyclisch (\cref{thm:b3:galois:cyclic}).
\end{solution}

\begin{solution}{exo:b3:representations:11}
(a) Equivariantie: $\rho(h)p_i\rho(h)^{-1}$ herindexeert de som
($g \mapsto hgh^{-1}$, en $\chi_i$ is een \omterm{def:b3:representations:character}{klassefunctie}): $p_i$
commuteert met de \omterm{def:b3:groups:action}{werking}. Op een irreducibele $W$ met \omterm{def:b3:representations:character}{karakter}
$\chi_j$ maakt Schur van de beperking een scalair
$\lambda\,\mathrm{id}_W$; sporen nemen geeft
\[
\lambda\,n_j = \frac{n_i}{\abs G}\sum_g
\overline{\chi_i(g)}\,\chi_j(g) = n_i\,\langle\chi_i,
\chi_j\rangle = n_i\,\delta_{ij}
\]
(de eerste orthogonaliteit): dus $\lambda = \delta_{ij}$.

(b) Ontbind $V$ in irreducibele (Maschke): $p_i$ werkt als de
identiteit op de summanden met \omterm{def:b3:representations:character}{karakter} $\chi_i$ en als $0$ op
alle andere, dus is $p_i$ de projectie op hun som $V_i$ langs de
som van de rest; het beeld $V_i$ hangt niet van de gekozen
ontbinding af (het is de verzameling vectoren die door $p_i$
worden vastgehouden, zonder keuzen gedefinieerd). $\sum_ip_i$
werkt als de identiteit op elke irreducibele summand: het is
$\mathrm{id}_V$. De fijnere splitsing van $V_i \cong
W_i^{\oplus m_i}$ vergt de keuze van een basis van
$\operatorname{Hom}_G(W_i, V)$: kanoniek is ze niet.

(c) Voor $\varepsilon$ (graad $1$) is $p_\varepsilon =
\frac1{6}\sum_{g}\varepsilon(g)\,\rho(g)$, dat wil zeggen in de
groepsalgebra
\[
p_\varepsilon = \tfrac16\bigl(e - (1\,2) - (1\,3) - (2\,3)
+ (1\,2\,3) + (1\,3\,2)\bigr) .
\]
Kwadrateren: de coëfficiënt van $g$ in $p_\varepsilon^2$ is
$\frac1{36}\sum_{h}\varepsilon(h)\varepsilon(h^{-1}g) =
\frac1{36}\,\varepsilon(g)\sum_h\varepsilon(h)^2 =
\frac{6}{36}\varepsilon(g)$: dus $p_\varepsilon^2 =
p_\varepsilon$. (Haar beeld in de \omterm{ex:b3:representations:examples}{reguliere representatie} is de
rechte opgespannen door $\sum_g\varepsilon(g)e_g$: de
tekenrepresentatie komt met multipliciteit $1$ voor, zoals de
algemene theorie eist.)
\end{solution}

\begin{solution}{exo:b3:representations:12}
(a) Orden de klassen als $e$ [1], de transposities [6], de
$3$-cykels [8], de $4$-cykels [6] en de dubbele transposities [3].
Het tweede lineaire \omterm{def:b3:representations:character}{karakter} is $\chi_2 = \varepsilon$ met
waarden $1, -1, 1, -1, 1$. Voor het \omterm{def:b3:representations:character}{karakter} $\chi_3$ van graad
$2$ geeft de \omterm{thm:b3:representations:orthogonality}{kolomorthogonaliteit} van elke kolom met die van de
identiteit ($\sum_in_i\chi_i(g) = 0$ voor $g \neq e$) op de
transposities: $1 - 1 + 2\chi_3 + 3(\chi_4 + \chi_5) = 0$; en de
tekentruc $\chi_5 = \varepsilon\chi_4$ (een \omterm{def:b3:representations:character}{karakter} van graad $3$
maal een lineair \omterm{def:b3:representations:character}{karakter} is opnieuw \omterm{def:b3:representations:rep}{irreducibel} --- dezelfde
norm) laat $\chi_4 + \chi_5$ op de oneven klassen verdwijnen: dus
$\chi_3 = 0$ daar. Op de $3$-cykels: $1 + 1 + 2\chi_3(c_3) +
3(\chi_4 + \chi_5)(c_3) = 0$ met $\chi_5 = \chi_4$ op de even
klassen; de kolom van $c_3$ met zichzelf levert gegevens over
$\abs{\chi_3}^2$; en het kleine stelsel oplossen (gebruik ook de
rijorthogonaliteit van $\chi_3$ met $\mathbf 1$ en $\varepsilon$)
geeft $\chi_3 = (2, 0, -1, 0, 2)$, en daarna $\chi_4 = (3, 1, 0,
-1, -1)$ en $\chi_5 = \varepsilon\chi_4 = (3, -1, 0, 1, -1)$. De
volledige tafel:
\begin{center}
\begin{tabular}{c|ccccc}
 & $e$ & $6\,t$ & $8\,c_3$ & $6\,c_4$ & $3\,v$\\
\hline
$\chi_1$ & $1$ & $1$ & $1$ & $1$ & $1$\\
$\chi_2$ & $1$ & $-1$ & $1$ & $-1$ & $1$\\
$\chi_3$ & $2$ & $0$ & $-1$ & $0$ & $2$\\
$\chi_4$ & $3$ & $1$ & $0$ & $-1$ & $-1$\\
$\chi_5$ & $3$ & $-1$ & $0$ & $1$ & $-1$\\
\end{tabular}
\end{center}
(Alle rijen hebben norm $1$ en alle kolommen zijn orthogonaal: de
controles slagen.)

(b) De klassegegevens $1, 6, 8, 6, 3$ met deze graden zijn die van
$S_4$ (cykeltypen $e$, $2$, $3$, $4$, $2{+}2$). Kernen:
$\ker\chi_2 = \{g : \varepsilon(g) = 1\} = A_4$ (de klassen $e,
c_3, v$: $1 + 8 + 3 = 12$, index $2$); $\ker\chi_3 = \{g :
\chi_3(g) = 2\}$ zijn de klassen $e$ en $v$: de viergroep van
Klein $V_4$, van orde $4$, normaal. En $\chi_4, \chi_5$ zijn trouw
($\chi_i(g) = n_i$ alleen in $e$). De doorsneden van kernen:
$\{e\}$, $V_4$, $A_4$, $G$ --- volgens
\cref{exo:b3:representations:6}(b) zijn dat \emph{alle}
\omterm{def:b3:groups:normal}{normaaldelers} van $S_4$.

(c) De \omterm{def:b3:representations:character}{karakters} met $V_4 \subseteq \ker$ zijn $\chi_1, \chi_2,
\chi_3$: ze factoriseren over $G/V_4$, een groep van orde $6$ met
irreducibele graden $1, 1, 2$ --- de tafel van $S_3$. Omdat de
tafel van het quotiënt onder de twee groepen van orde $6$ een
volledige invariant \emph{is} ($\Z/6\Z$ zou zes lineaire \omterm{def:b3:representations:character}{karakters}
hebben), is $G/V_4 \cong S_3$: het quotiënt is in de tafel
zichtbaar als het blok rijen dat $V_4$ in zijn kern bevat.
\end{solution}

\begin{solution}{pb:b3:representations:1}
\textbf{1.} Is $\alpha^n + c_{n-1}\alpha^{n-1} + \dots + c_0 = 0$
met $c_i \in \Z$, dan ligt $\alpha^n$ in het $\Z$-opspansel van
$1, \dots, \alpha^{n-1}$, en inductief geldt dat voor elke macht:
$\Z[\alpha]$ wordt voortgebracht door $1, \alpha, \dots,
\alpha^{n-1}$. Omgekeerd, zij $\Z[\alpha] = \Z g_1 + \dots + \Z
g_m$. Schrijf $\alpha g_i = \sum_j m_{ij}g_j$ met $M = (m_{ij})
\in M_m(\Z)$: de vector $g = (g_i)$ voldoet aan $(\alpha I - M)g =
0$; vermenigvuldigen met de geadjugeerde matrix geeft $\det(\alpha
I - M)\,g_i = 0$ voor elke $i$, en omdat $1 \in \Z[\alpha]$ een
$\Z$-combinatie van de $g_i$ is, volgt $\det(\alpha I - M) = 0$:
$\alpha$ is een wortel van de monische $\det(XI - M) \in \Z[X]$.

\textbf{2.} Wordt $\Z[\alpha]$ opgespannen door de machten van
$\alpha$ tot en met $n-1$ en $\Z[\beta]$ door die van $\beta$ tot
en met $m - 1$, dan wordt $\Z[\alpha,\beta]$ opgespannen door de
$nm$ producten $\alpha^i\beta^j$ (reduceer elk monoom). De
deelringen $\Z[\alpha + \beta]$ en $\Z[\alpha\beta]$ zijn
$\Z$-deelmodulen van het eindig voortgebrachte $\Z$-moduul
$\Z[\alpha, \beta]$ en dus \omterm{def:b3:modules:free}{eindig voortgebracht} ($\Z[\alpha,
\beta]$, voortgebracht door $nm$ elementen, is een beeld van
$\Z^{nm}$; een deelmoduul trekt terug tot een deelmoduul van
$\Z^{nm}$, vrij van rang $\leq nm$ volgens
\cref{thm:b3:modules:submodule}, en het beeld daarvan brengt
voort). Volgens vraag 1 zijn $\alpha + \beta$ en $\alpha\beta$
gehele \omterm{def:b3:galois:algebraic}{algebraïsche} getallen.

\textbf{3.} Is $\frac pq$ (in laagste termen) een wortel van een
monische gehele veelterm van graad $n$, dan geeft de
rationalewortelstelling (werk de noemers weg: $p^n = -q(\cdots)$)
dat $q \mid p^n$, dus $q = \pm1$.

\textbf{4.} $\chi(g)$ is een som van eenheidswortels
(\cref{prop:b3:representations:charbasics}), elk een geheel
\omterm{def:b3:galois:algebraic}{algebraïsch} getal (wortel van $X^N - 1$); besluit met vraag 2.

\textbf{5.} Voor $h \in G$ is $\rho(h)S_C\rho(h)^{-1} = \sum_{g\in
C}\rho(hgh^{-1}) = S_C$ (want $C$ is een klasse). Volgens Schur is
$S_C = \omega_C\,\mathrm{id}$; sporen nemen geeft $\abs
C\,\chi(g_C) = \omega_C\, n$.

\textbf{6.} $S_CS_{C'} = \sum_{x \in C, y \in C'}\rho(xy) =
\sum_{z \in G} a(z)\rho(z)$ met $a(z) = \#\{(x,y) \in C\times C' :
xy = z\}$. Conjugatie met $h$ brengt de oplossingen voor $z$
bijectief naar die voor $hzh^{-1}$: $a$ is dus een \omterm{def:b3:representations:character}{klassefunctie}
met waarden in $\N$, zodat $S_CS_{C'} =
\sum_{C''}a_{CC'C''}S_{C''}$. Overal $S_C = \omega_C\,\mathrm{id}$
substitueren en de scalairen identificeren geeft
$\omega_C\omega_{C'} = \sum_{C''}a_{CC'C''}\,\omega_{C''}$.

\textbf{7.} Zij $M$ het $\Z$-moduul opgespannen door $1$ en alle
producten $\omega_{C_1}\cdots\omega_{C_k}$; volgens vraag 6
reduceert elk zulk product tot een $\Z$-combinatie van $1$ en de
$\omega_C$: dus is $M$ \omterm{def:b3:modules:free}{eindig voortgebracht}, en $\omega_C M
\subseteq M$ voor elke $C$. In het bijzonder is $\Z[\omega_C]
\subseteq M$ \omterm{def:b3:modules:free}{eindig voortgebracht} (deelmoduul, als in vraag 2), en
maakt vraag 1 van $\omega_C$ een geheel \omterm{def:b3:galois:algebraic}{algebraïsch} getal.

\textbf{8.} Voor een irreducibele $\chi$ van graad $n$:
\[
\frac{\abs G}{n} = \frac{\abs G}{n}\langle\chi,\chi\rangle
= \frac1n \sum_{g}\chi(g)\overline{\chi(g)}
= \sum_{C}\frac{\abs C\,\chi(g_C)}{n}\,\overline{\chi(g_C)}
= \sum_C \omega_C\,\overline{\chi(g_C)} .
\]
Elke $\overline{\chi(g_C)} = \chi(g_C^{-1})$ is een geheel
\omterm{def:b3:galois:algebraic}{algebraïsch} getal (vraag 4), dus is het rechterlid er ook een
(vragen 2 en 7); en het is rationaal, dus een geheel getal (vraag
3): $n \mid \abs G$.

\textbf{9.} Bézout: $u\abs C + vn = 1$ met $u, v \in \Z$. Dan is
\[
\frac{\chi(g_C)}{n} = u\,\frac{\abs C\,\chi(g_C)}{n} +
v\,\chi(g_C) = u\,\omega_C + v\,\chi(g_C),
\]
een geheel \omterm{def:b3:galois:algebraic}{algebraïsch} getal.

\textbf{10.} Stel $0 < \abs{\chi(g_C)} < n$ en zij $\alpha =
\chi(g_C)/n$. Alle waarden $\chi(g)$ liggen in $\Q(\zeta_N)$ met
$N = \abs G$ (sommen van $N$-de eenheidswortels). Voor $\sigma \in
\operatorname{Gal}(\Q(\zeta_N)/\Q)$ stuurt $\sigma$
eenheidswortels naar eenheidswortels ($\sigma(\zeta^k) =
\zeta^{ak}$, \cref{thm:b3:galois:cyclotomicirred}), zodat
$\sigma(\chi(g_C))$ opnieuw een som van $n$ eenheidswortels is:
$\abs{\sigma(\alpha)} \leq 1$; en $\sigma(\alpha)$ is een geheel
\omterm{def:b3:galois:algebraic}{algebraïsch} getal (het heeft dezelfde \omterm{def:b3:galois:algebraic}{minimaalveelterm} als
$\alpha$). Het product $P = \prod_\sigma \sigma(\alpha)$ wordt
door de hele \omterm{def:b3:galois:galois}{Galoisgroep} vastgehouden en is dus rationaal
(\cref{thm:b3:galois:fundamental}(1)), en het is een geheel
\omterm{def:b3:galois:algebraic}{algebraïsch} getal met
\[
0 < \abs P \leq \abs\alpha < 1
\]
(geen enkele factor verdwijnt: $\sigma(\alpha) = 0$ zou $\alpha =
0$ afdwingen). Dat spreekt vraag 3 tegen. Bijgevolg is $\chi(g_C)
= 0$ of $\abs{\chi(g_C)} = n$, en in het laatste geval is
$\rho(g_C)$ scalair
(\cref{prop:b3:representations:charbasics}).

\textbf{11.} \omterm{thm:b3:representations:orthogonality}{Kolomorthogonaliteit} ($C \neq \{e\}$): $\sum_i
\chi_i(e)\overline{\chi_i(g_C)} = 0$, dat wil zeggen $1 + \sum_{i
\geq 2} n_i\overline{\chi_i(g_C)} = 0$. Verdween elk
niet-triviaal $\chi_i$ met $p \nmid n_i$ in $g_C$, dan gaf
groeperen van de rest naar hun factor $p$:
\[
-\frac1p = \sum_{i \geq 2,\ p \mid n_i}
\frac{n_i}{p}\,\overline{\chi_i(g_C)},
\]
een geheel \omterm{def:b3:galois:algebraic}{algebraïsch} getal --- in strijd met vraag 3. Er is dus
een niet-triviaal $\chi_i$ met $p \nmid n_i$ en $\chi_i(g_C) \neq
0$; en omdat $\abs C = p^k$ is $\gcd(\abs C, n_i) = 1$, zodat
vraag 10 van $\rho_i(g_C)$ een scalair maakt. Nu is $G$
enkelvoudig en niet-abels: $\chi_i$ is trouw
(\cref{exo:b3:representations:6}(c)), en $Z_i = \{g : \rho_i(g)
\text{ scalair}\}$ is een \omterm{def:b3:groups:normal}{normaaldeler} (het origineel onder
$\rho_i$ van de scalairen, die een normale --- zelfs centrale ---
deelgroep van het beeld vormen) die $g_C \neq e$ bevat: dus $Z_i =
G$. Dan is $\rho_i(G)$ abels en trouw, zodat $G$ abels is:
tegenspraak. \emph{Geen enkele enkelvoudige niet-abelse groep
heeft een conjugatieklasse van priemmachtgrootte $> 1$.}

\textbf{12.} Zij $G$ enkelvoudig van orde $p^aq^b$. Is $b = 0$
($G$ een $p$-groep), dan is $Z(G) \neq \{e\}$
(\cref{thm:b3:groups:pfixed}) normaal, dus $Z(G) = G$: abels.
Neem anders een Sylow-$q$-deelgroep $Q$ en $z \in Z(Q)
\setminus\{e\}$ (opnieuw \cref{thm:b3:groups:pfixed}). Dan is $Q
\subseteq Z_G(z)$, zodat de klasse van $z$ grootte $[G : Z_G(z)]$
heeft, een deler van $[G : Q] = p^a$: een macht van $p$. Is die
grootte $1$, dan is $z \in Z(G)$: het \omterm{ex:b3:groups:actions}{centrum} is een
niet-triviale \omterm{def:b3:groups:normal}{normaaldeler}, dus $Z(G) = G$, abels. Is ze $p^k >
1$, dan verbiedt vraag 11 dat voor een enkelvoudige niet-abelse
$G$. Hoe dan ook is een \omterm{def:b3:groups:simple}{enkelvoudige groep} van orde $p^aq^b$
abels ($\cong \Z/p\Z$).

\textbf{13.} Inductie naar $\abs G$ (voor $\abs G = 1$:
\omterm{def:b3:groups:derived}{oplosbaar}). Is $G$ enkelvoudig, dan maakt vraag 12 haar abels en
dus \omterm{def:b3:groups:derived}{oplosbaar}. Kies anders een niet-triviale echte \omterm{def:b3:groups:normal}{normaaldeler}
$N$: $\abs N$ en $\abs{G/N}$ zijn opnieuw van de vorm
$p^{a'}q^{b'}$ en kleiner, dus zijn $N$ en $G/N$ per inductie
\omterm{def:b3:groups:derived}{oplosbaar} en is $G$ \omterm{def:b3:groups:derived}{oplosbaar}
(\cref{prop:b3:groups:derived}). --- Bij drie priemgetallen faalt
de sleutelstap: de index van een \omterm{def:b3:groups:sylow}{Sylow-deelgroep} is dan geen
priemmacht meer, zodat de klasse van een centraal element van een
\omterm{def:b3:groups:sylow}{Sylow-deelgroep} geen priemmachtgrootte hoeft te hebben. En
falen moet het wel: $A_5$, van orde $2^2\cdot3\cdot5$, is
enkelvoudig en niet \omterm{def:b3:groups:derived}{oplosbaar}.

\textbf{14.} De klassen en groottes zijn die van
\cref{exo:b3:groups:11}(a). De $S_5$-klasse van $c$ heeft grootte
$24$; bleef ze één $A_5$-klasse, dan zou de telling met \omterm{def:b3:groups:action}{baan} en
\omterm{def:b3:groups:action}{stabilisator} $\abs{Z_{A_5}(c)} = 60/24$ geven, geen geheel getal
--- concreet heeft $Z_{S_5}(c) = \langle c\rangle$ orde $5$ en
ligt die binnen $A_5$, zodat de $A_5$-klasse van $c$ grootte
$60/5 = 12$ heeft: de $S_5$-klasse valt in twee uiteen ($c$ en
$c^2$ zijn in $S_5$ alleen door een oneven permutatie
geconjugeerd).

\textbf{15.} Eén lineair \omterm{def:b3:representations:character}{karakter}: een \omterm{def:b3:representations:rep}{representatie} van graad $1$
factoriseert over $G/D(G)$, en $D(A_5) = A_5$ ($A_5$ is
enkelvoudig en niet-abels, en $D(A_5)$ is normaal en
niet-triviaal). Dus $n_1 = 1$ en $n_2^2 + n_3^2 + n_4^2 + n_5^2 =
59$ met elke $n_i \geq 2$. Beschikbare kwadraten: $4, 9, 16, 25,
36, 49$. Een som van vier ervan gelijk aan $59$: het grootste moet
$25$ zijn (met $36$ of $49$ lukt geen enkele combinatie: $36 + 16
+ 4 + 4 = 60$, $36 + 9 + 9 + 4 = 58$, $36 + 16 + 9 + 4 = 65$), en
$59 - 25 = 34 = 16 + 9 + 9$ (de enige manier: $16 + 16 + 4 = 36$,
$25 + 9 + 4 = 38$, $25 + 4 + 4 = 33$): graden $1, 3, 3, 4, 5$.

\textbf{16.} $\langle\pi, \mathbf 1\rangle = \frac1{60}(5 +
15\cdot1 + 20\cdot2 + 0 + 0) = 1$ (één \omterm{def:b3:groups:action}{baan} --- de telling van
Burnside), en $\langle\pi, \pi\rangle = \frac1{60}(25 + 15 + 80 +
0 + 0) = 2$: $\pi$ bevat het triviale \omterm{def:b3:representations:character}{karakter} één keer, en haar
andere bestanddeel is één enkel \omterm{def:b3:representations:rep}{irreducibel} \omterm{def:b3:representations:character}{karakter}. Bijgevolg is
$\chi_4 = \pi - \mathbf 1$ \omterm{def:b3:representations:rep}{irreducibel} van graad $4$, met waarden
$4, 0, 1, -1, -1$.

\textbf{17.} Op de paren houdt een element $\{i,j\}$ vast precies
wanneer het $i$ en $j$ vasthoudt of verwisselt: de aantallen zijn
$10$ (voor $e$), $2$ ($t = (1\,2)(3\,4)$ houdt $\{1,2\}$ en
$\{3,4\}$ vast), $1$ ($(1\,2\,3)$ houdt $\{4,5\}$ vast), $0$ en
$0$. Dan is $\langle\pi_{10}, \pi_{10}\rangle = \frac1{60}(100 +
15\cdot4 + 20\cdot1) = 3$: drie irreducibele bestanddelen, elk één
keer. En $\langle\pi_{10}, \mathbf 1\rangle = \frac1{60}(10 + 30 +
20) = 1$ en $\langle\pi_{10}, \chi_4\rangle = \frac1{60}(40 + 0 +
20\cdot1\cdot1 + 0 + 0) = 1$: dus $\pi_{10} = \mathbf 1 + \chi_4 +
\chi$ met $\chi$ \omterm{def:b3:representations:rep}{irreducibel} van graad $10 - 1 - 4 = 5$ en waarden
$\chi_5 = \pi_{10} - \mathbf 1 - \chi_4 = (5, 1, -1, 0, 0)$.

\textbf{18.} Schrijf $a, a'$ voor de waarden van $\chi_2, \chi_3$
in $t$, en $b, b'$ voor die in $s = (1\,2\,3)$; alle vier zijn
reëel ($t$ en $s$ zijn met hun inversen geconjugeerd). Kolom $t$
tegen kolom $e$: $1 + 3a + 3a' + 4\cdot0 + 5\cdot1 = 0$, dus $a +
a' = -2$; kolom $t$ met zichzelf: $1 + a^2 + a'^2 + 0 + 1 =
\frac{60}{15} = 4$, dus $a^2 + a'^2 = 2$; en dan geeft $(a + a')^2
= 4 = 2 + 2aa'$ dat $aa' = 1$ en $a = a' = -1$. Kolom $s$ tegen
$e$: $1 + 3(b + b') + 4 - 5 = 0$ geeft $b + b' = 0$; en kolom $s$
met zichzelf: $1 + b^2 + b'^2 + 1 + 1 = \frac{60}{20} = 3$ geeft
$b = b' = 0$.

\textbf{19.} Kolom $c$ tegen kolom $e$: $1 + 3(x + y) + 4(-1) +
5\cdot0 = 0$, dus $x + y = 1$. Kolom $c$ tegen kolom $c^2$
(verschillende klassen, dus orthogonaal): $1 + xy + yx + 1 + 0 =
0$, dus $xy = -1$. Bijgevolg lossen $x$ en $y$ de vergelijking
$T^2 - T - 1 = 0$ op: $\{x, y\} = \{\varphi, \bar\varphi\}$ met
$\varphi = \frac{1 + \sqrt5}2$. Consistentie: $x^2 + y^2 = (x+y)^2
- 2xy = 3 = \frac{60}{12} - 2$ --- in overeenstemming met de
identiteit voor de kolom met zichzelf, $1 + x^2 + y^2 + 1 + 0 =
5$. De tafel:
\begin{center}
\begin{tabular}{c|ccccc}
 & $e$ & $15\,t$ & $20\,s$ & $12\,c$ & $12\,c^2$\\
\hline
$\chi_1$ & $1$ & $1$ & $1$ & $1$ & $1$\\
$\chi_2$ & $3$ & $-1$ & $0$ & $\varphi$ & $\bar\varphi$\\
$\chi_3$ & $3$ & $-1$ & $0$ & $\bar\varphi$ & $\varphi$\\
$\chi_4$ & $4$ & $0$ & $1$ & $-1$ & $-1$\\
$\chi_5$ & $5$ & $1$ & $-1$ & $0$ & $0$\\
\end{tabular}
\end{center}

\textbf{20.} $\norm{\chi_2}^2 = \frac1{60}\bigl(9 + 15\cdot1 + 0 +
12\varphi^2 + 12\bar\varphi^2\bigr) = \frac{9 + 15 + 36}{60} = 1$,
met $\varphi^2 + \bar\varphi^2 = (\varphi + \bar\varphi)^2 -
2\varphi\bar\varphi = 1 + 2 = 3$. Net zo is $\langle\chi_2,
\chi_3\rangle = \frac1{60}(9 + 15 + 0 + 12(2\varphi\bar\varphi)) =
\frac{9 + 15 - 24}{60} = 0$. Deelbaarheid van \omterm{thm:b3:galois:finitefields}{Frobenius}: $1, 3, 3,
4, 5$ delen alle $60$. De irrationale waarden $\varphi$ en
$\bar\varphi$ zijn het zichtbare verschil met $S_5$: in $S_5$ is
elk element geconjugeerd met alle voortbrengers van zijn cyclische
groep met hetzelfde cykeltype --- in het bijzonder $c \sim c^2$
--- wat rationale (zelfs gehele) karakterwaarden afdwingt; in
$A_5$ opent de splitsing van de $5$-cykels de deur naar
$\Q(\sqrt5)$.

\textbf{21.} Een \omterm{def:b3:groups:normal}{normaaldeler} $N$ is een vereniging van klassen,
bevat $e$, en $\abs N \mid 60$. De kandidaatsommen: $1{+}15 = 16$,
$1{+}20 = 21$, $1{+}12 = 13$, $1{+}24 = 25$, $1{+}15{+}20 = 36$,
$1{+}15{+}12 = 28$, $1{+}15{+}24 = 40$, $1{+}20{+}12 = 33$,
$1{+}20{+}24 = 45$, $1{+}12{+}12 = 25$, $1{+}15{+}20{+}12 = 48$,
en zo alle echte deelsommen die $1$ bevatten: geen van $13, 16,
21, 25, 28, 33, 36, 40, 45, 48$ deelt $60$, en dus blijft alleen
$N = \{e\}$ of $N = A_5$ over. Enkelvoudigheid, afgelezen uit vijf
getallen. Ook het criterium van vraag 11 is zichtbaar: de
klassegroottes $15 = 3\cdot5$, $20 = 4\cdot5$ en $12 = 4\cdot3$
zijn alle samengesteld uit twee priemgetallen --- geen enkele
klasse van priemmachtgrootte, precies zoals het criterium van
Burnside van een \omterm{def:b3:groups:simple}{enkelvoudige groep} eist.

\textbf{22.} Een rotatie van $\R^3$ over een hoek $\theta$ heeft
eigenwaarden $1, \eu^{\iu\theta}, \eu^{-\iu\theta}$: spoor $1 +
2\cos\theta$. Voor $\theta = \frac{2\pi}5$ is $2\cos\frac{2\pi}5 =
\frac{\sqrt5 - 1}2$, dus is het spoor $1 + \frac{\sqrt5-1}2 =
\frac{1+\sqrt5}2 = \varphi$. Het niet-triviale element $\tau$ van
$\operatorname{Gal}(\Q(\sqrt5)/\Q)$, ingang voor ingang op een
\omterm{def:b3:representations:table}{karaktertafel} toegepast, stuurt \omterm{def:b3:representations:character}{karakters} naar \omterm{def:b3:representations:character}{karakters} (het
verdraagt zich met de definiërende algebra: $\tau\circ\chi$ is het
\omterm{def:b3:representations:character}{karakter} van de \omterm{def:b3:representations:rep}{representatie} die je krijgt door de matrices via
$\tau$ op de ingangen over te brengen, of abstract: de
\omterm{thm:b3:representations:orthogonality}{orthogonaliteitsrelaties} zijn $\Q$-rationaal, dus permuteert
$\tau$ hun oplossingen); $\tau$ houdt $\chi_1, \chi_4, \chi_5$
vast (rationale waarden) en moet dus $\chi_2$ en $\chi_3$
verwisselen: het tweede \omterm{def:b3:representations:character}{karakter} van graad $3$ draagt de
toegevoegde waarden, zonder dat er één matrix is berekend.
Meetkundig zijn de twee \omterm{def:b3:representations:rep}{representaties} de icosaëdrische \omterm{def:b3:groups:action}{werking} en
haar samenstelling met een uitwendig automorfisme van $A_5$
(conjugatie met een transpositie), dat de twee klassen van
$5$-cykels verwisselt.

\textbf{23.} Voor $z \in Z(G)$ commuteert $\rho(z)$ met elke
$\rho(g)$, dus is volgens het lemma van Schur $\rho(z) =
\lambda\,\mathrm{id}$; en omdat $z$ eindige orde heeft, is
$\lambda$ een eenheidswortel en $\abs{\chi(z)} = \abs\lambda\,n =
n$. Dan is
\[
\abs G = \abs G\,\langle\chi, \chi\rangle
= \sum_{g \in G}\abs{\chi(g)}^2
\geq \sum_{z \in Z(G)}\abs{\chi(z)}^2
= \abs{Z(G)}\,n^2,
\]
dat wil zeggen $n^2 \leq [G : Z(G)]$. Een niet-abelse groep van
orde $8$ ($D_4$ of $Q_8$) heeft vijf \omterm{ex:b3:groups:actions}{conjugatieklassen}, dus vijf
irreducibele graden met $\sum n_i^2 = 8$; en de enige manier om
$8$ als som van vijf kwadraten $\geq 1$ te schrijven is $1 + 1 + 1
+ 1 + 4$: graden $1, 1, 1, 1, 2$. Beide groepen hebben een \omterm{ex:b3:groups:actions}{centrum}
van orde $2$, en het \omterm{def:b3:representations:character}{karakter} van graad $2$ bereikt de gelijkheid:
$2^2 = 4 = [G : Z(G)]$ --- de grens is scherp. Voor $A_5$ is $Z =
\{e\}$ en luidt de grens $n^2 \leq 60$: ruimschoots voldaan door
$1, 3, 3, 4, 5$ ($25 \leq 60$), zoals het moet, want gelijkheid
zou (langs dezelfde keten) afdwingen dat $\chi$ buiten het \omterm{ex:b3:groups:actions}{centrum}
verdwijnt.

\textbf{24.} Er geldt $\rho(g)^m = \mathrm{id}$ met $m$ de orde
van $g$, dus wordt $\rho(g)$ geannihileerd door $X^m - 1$, die
over $\C$ uiteenvalt met enkelvoudige wortels: $\rho(g)$ is
diagonaliseerbaar met eenheidswortels $\lambda_1, \dots,
\lambda_n$ als eigenwaarden, in een eigenbasis $(e_i)$. De
producten $e_ie_j$ ($i \leq j$) vormen een eigenbasis van
$\operatorname{Sym}^2V$ met eigenwaarden $\lambda_i\lambda_j$, en
de $e_i \wedge e_j$ ($i < j$) een van $\Lambda^2V$; omdat
\[
\sum_{i<j}\lambda_i\lambda_j
= \frac{(\sum_i\lambda_i)^2 - \sum_i\lambda_i^2}2
= \frac{\chi(g)^2 - \chi(g^2)}2,
\]
en de symmetrische som $\sum_i\lambda_i^2$ optelt in plaats van
aftrekt, volgen beide formules. Voor $\chi_2 = (3, -1, 0, \varphi,
\bar\varphi)$ op de klassen $(e, (2,2)\text{-},
3\text{-cykels}, c, c^2)$: kwadrateren stuurt de dubbele
transposities naar $e$, de drietallige cykels naar drietallige
cykels, en de klasse van $c$ op die van $c^2$ en omgekeerd (in
$A_5$ is $c^{-1} \sim c$, dus $c^4 \sim c$). Bijgevolg leest
$\chi_2(g^2)$ als $(3, 3, 0, \bar\varphi, \varphi)$, en met
$\varphi^2 = \varphi + 1$ en $\bar\varphi = 1 - \varphi$:
\[
\Lambda^2\chi_2 = (3, -1, 0, \varphi, \bar\varphi) = \chi_2,
\qquad
\operatorname{Sym}^2\chi_2 = (6, 2, 0, 1, 1) .
\]
Dat laatste ontbinden geeft, met de klassegroottes $1, 15, 20, 12,
12$: $\langle\cdot, \mathbf 1\rangle = \frac1{60}(6 + 15\cdot2 + 0
+ 12 + 12) = 1$; $\langle\cdot, \chi_5\rangle = \frac1{60}(30 +
30) = 1$; $\langle\cdot, \chi_4\rangle = \frac1{60}(24 - 12 - 12)
= 0$; en $\langle\cdot, \chi_2\rangle = \frac1{60}(18 - 30 +
12(\varphi + \bar\varphi)) = 0$, en net zo voor $\chi_3$. Dus is
$\operatorname{Sym}^2\chi_2 = \mathbf 1 + \chi_5$ (dimensies $6 =
1 + 5$) en $\chi_2\otimes\chi_2 = \mathbf 1 + \chi_2 + \chi_5$
(dimensies $9 = 1 + 3 + 5$). Meetkundig: het equivariante
isomorfisme $\Lambda^2\R^3 \to \R^3$, $u \wedge v \mapsto u \times
v$, is precies $\Lambda^2\chi_2 = \chi_2$ voor een rotatiegroep;
de summand $\mathbf 1$ van het symmetrische kwadraat is de
invariante kwadratische vorm $x^2 + y^2 + z^2$, en $\chi_5$ leeft
op de vijfdimensionale ruimte van spoorloze symmetrische tensoren
(de harmonische kwadratische vormen).

\textbf{25.} Kolom van $c$ tegen kolom van $c^2$:
\[
1\cdot1 + \varphi\bar\varphi + \bar\varphi\varphi +
(-1)(-1) + 0 = 1 - 1 - 1 + 1 + 0 = 0,
\]
zoals de orthogonaliteit voor verschillende klassen eist (want
$\varphi \bar\varphi = -1$). Kolom van $c$ tegen zichzelf: $1 +
\varphi^2 + \bar\varphi^2 + 1 + 0 = 1 + 3 + 1 = 5 = 60/12 =
\abs{Z_{A_5}(c)}$. Het reguliere \omterm{def:b3:representations:character}{karakter} $\sum_in_i\chi_i$ op de
vier kolommen die niet bij de identiteit horen:
\[
1 - 3 - 3 + 0 + 5 = 0, \qquad
1 + 0 + 0 + 4 - 5 = 0,
\]
op de dubbele transposities en de drietallige cykels, en op de
klasse van $c$ (die van $c^2$ is haar Galois-toegevoegde):
\[
1 + 3\varphi + 3\bar\varphi - 4 + 0 = 1 + 3 - 4 = 0,
\]
met $\varphi + \bar\varphi = 1$. De tafel doorstaat elke
controle: het is de \omterm{def:b3:representations:table}{karaktertafel} van $A_5$.
\end{solution}
